% ============================================================
\chapter{Modul 2: Entscheidungen treffen}
% ============================================================

Ein Roboter, der nur stur geradeaus fährt, ist kein autonomes System -- er ist ein ferngesteuertes Fahrzeug. Autonomie beginnt mit der Fähigkeit, Situationen zu beurteilen und zu reagieren. Genau das lernen wir hier.

\section{Lektion 5: Bedingungen mit \texttt{if}}

\begin{analogie}
Jeder Fahrstuhl kennt folgende Logik: \textit{Wenn Tür offen, dann fahre nicht.} Das ist eine Bedingung. Bedingungen sind das Grundprinzip aller Sicherheitssysteme -- und aller autonomen Roboter.
\end{analogie}

\begin{lstlisting}
abstand = 0.18   # Meter, gemessen vom Ultraschallsensor

if abstand < 0.20:
    print("Hindernis erkannt!")
    print("Stoppe Motoren.")
\end{lstlisting}

\begin{merksatz}
Die Syntax: \texttt{if BEDINGUNG:} -- danach eingerückter Block. Einrückung ist in Python \textbf{Pflicht}, nicht Stil. Falsche Einrückung = Fehler.
\end{merksatz}

\subsection*{if -- elif -- else}

\begin{lstlisting}
abstand = 0.35   # Meter

if abstand < 0.20:
    print("STOP: Hindernis zu nah!")
elif abstand < 0.50:
    print("WARNUNG: Langsam fahren.")
elif abstand < 1.00:
    print("OK: Normal fahren.")
else:
    print("FREI: Volle Geschwindigkeit.")
\end{lstlisting}

\begin{praxis}
Dieses Muster -- Abstandsstufen mit unterschiedlichen Reaktionen -- ist der Kern jeder Hindernisvermeidung. In echten Systemen heißen die Schwellenwerte Parameter und werden aus Konfigurationsdateien geladen (siehe Modul 7).
\end{praxis}

\begin{uebung}[title=Übung 2.1: Akkumanagement]
Schreibe eine Bedingungslogik für den Akkustand:
\begin{itemize}
    \item unter 10\%: ``KRITISCH -- Sofort laden!''
    \item unter 25\%: ``NIEDRIG -- Ladestation anfahren.''
    \item unter 50\%: ``MITTEL -- Aufgabe beenden, dann laden.''
    \item sonst: ``GUT -- Weiter arbeiten.''
\end{itemize}
\end{uebung}

\section{Lektion 6: Vergleichsoperatoren}

\begin{center}
\begin{tabular}{lll}
\toprule
\textbf{Operator} & \textbf{Bedeutung} & \textbf{Robotik-Beispiel} \\
\midrule
\texttt{==} & gleich            & \texttt{status == "bereit"} \\
\texttt{!=} & ungleich          & \texttt{fehler != None} \\
\texttt{<}  & kleiner           & \texttt{akku < 20} \\
\texttt{>}  & größer            & \texttt{temp > 80} \\
\texttt{<=} & kleiner gleich    & \texttt{last < max\_last} \\
\texttt{>=} & größer gleich     & \texttt{spannung >= 11.0} \\
\bottomrule
\end{tabular}
\end{center}

\begin{merksatz}
\texttt{=} weist einen Wert zu. \texttt{==} vergleicht. Dieser Unterschied ist eine der häufigsten Fehlerquellen in Python. \texttt{akku = 20} speichert 20. \texttt{akku == 20} fragt: \textit{Ist akku gleich 20?}
\end{merksatz}

\section{Lektion 7: Logische Operatoren}

\begin{lstlisting}
spannung     = 12.1
strom        = 0.45
notaus       = False
sensor_ok    = True

# and: beide muessen True sein
if spannung >= 11.0 and strom < 2.0:
    print("Elektrische Parameter: OK")

# or: mindestens eine muss True sein
if spannung < 10.0 or notaus:
    print("STOP: Sicherheitsunterbrechung!")

# not: kehrt True/False um
if not sensor_ok:
    print("WARNUNG: Sensor nicht erreichbar.")

# Kombination: Prioritaetslogik
if notaus:
    print("NOTAUS -- Alles stoppen.")
elif spannung < 11.0 or strom > 3.0:
    print("Elektrisches Problem erkannt.")
elif sensor_ok and spannung >= 11.0:
    print("System betriebsbereit.")
\end{lstlisting}

\begin{praxis}
Reale Sicherheitssysteme kombinieren immer mehrere Bedingungen. Ein Roboterarm darf sich nur bewegen, wenn: Tür geschlossen \texttt{and} Notaus nicht aktiv \texttt{and} Stromversorgung stabil. Alle drei müssen gleichzeitig \texttt{True} sein.
\end{praxis}

\begin{uebung}[title=Projektaufgabe Modul 2: Sicherheits-Check]
Schreibe einen vollständigen Sicherheits-Check für den Roboter-Start. Prüfe mindestens vier Bedingungen (Akku, Sensoren, Notaus, Kommunikation). Nutze \texttt{and}, \texttt{or}, \texttt{not}. Gib klare Statusmeldungen aus.
\end{uebung}

\begin{ros}
In ROS2 entspricht diese Logik dem \textit{Lifecycle Node}: Ein Node geht nur in den Zustand ``aktiv'', wenn alle Vorbedingungen erfüllt sind. Die Prüflogik sieht fast identisch aus wie der Code in dieser Übung.
\end{ros}
